% Generated draft: requires mathematical and source review.
\documentclass[11pt]{article}
\usepackage[margin=25mm]{geometry}
\usepackage{fontspec}
\setmainfont{Latin Modern Roman}
\usepackage{amsmath,amssymb,mathtools,unicode-math}
\setmathfont{Latin Modern Math}
\usepackage{enumitem,needspace,xurl,hyperref}
\setlength{\parindent}{0pt}
\setlength{\parskip}{5pt}
\setlength{\emergencystretch}{3em}
\newcommand{\N}{\mathbb{N}}
\newcommand{\Z}{\mathbb{Z}}
\newcommand{\Q}{\mathbb{Q}}
\newcommand{\R}{\mathbb{R}}
\newcommand{\C}{\mathbb{C}}
\newcommand{\F}{\mathbb{F}}
\newcommand{\A}{\mathbb{A}}
\newcommand{\PP}{\mathbb{P}}
\newcommand{\E}{\mathbb{E}}
\newcommand{\eps}{\varepsilon}
\newcommand{\dd}{\,\mathrm{d}}
\begin{document}
\section*{UnsolvedMath ID: 30004341}
\textbf{Source problem code:} catalogue entry
\subsection{Definitions and mathematical statement}
Let \(\overline{\Q}\subset\C\) denote the field of algebraic numbers. A singular modulus is a value \(j(\tau)\), where \(j\) is the classical modular invariant and \(\tau\in\C\) satisfies \(\operatorname{Im}(\tau)>0\) and \([\Q(\tau):\Q]=2\). Write \(\mathcal J\subset\overline{\Q}\) for the set of all singular moduli.

For a finite set \(S\) of rational primes, call \(u\in\overline{\Q}^{\times}\) an \(S\)-unit if, in a number field \(K\) containing \(u\), one has
\[
 v_{\mathfrak p}(u)=0
 \quad\text{for every nonzero prime ideal }\mathfrak p\subset\mathcal O_K
 \text{ lying above a rational prime outside }S.
\]
Here \(v_{\mathfrak p}\) is the additive prime-ideal valuation. This condition is independent of the choice of \(K\). Denote the set of these algebraic \(S\)-units by \(U_S\).

\textbf{Question (under the stated interpretation).} For every finite set \(S\) of rational primes, is the set
\[
 D_S=\{j_1-j_2:\ j_1,j_2\in\mathcal J,\ j_1\ne j_2,\ j_1-j_2\in U_S\}
\]
finite?

Both singular moduli vary. The requested finiteness concerns distinct difference values; the supplied wording does not explicitly ask for finiteness of the pairs producing them. Zero is excluded because an \(S\)-unit is nonzero. No bound independent of \(S\) is asserted.

The original question says only a fixed set of primes. The finiteness of \(S\) is taken from the supplied literature assessment, and the rational-prime and algebraic \(S\)-unit conventions above are explicit interpretations requiring source confirmation. The supplied record labels the problem open. Its dated assessment states that finiteness is known when one singular modulus is fixed, whereas finiteness as both vary remains conjectural; these are supplied curation claims, not independently verified mathematical or literature conclusions.
\subsection{Short English statement}
For each fixed finite set \(S\) of rational primes, ask whether only finitely many distinct nonzero differences of two singular moduli are algebraic \(S\)-units. Both moduli may vary. The finite-set hypothesis and the meaning of \(S\)-unit are made explicit from the supplied context and require human source review.
\subsection{Sources}
\begin{itemize}
\item Local UnsolvedMath registry, \texttt{data/unsolvedmath/problems.json}: permanent ID \(30004341\), title \emph{S-Unit Differences of Singular Moduli}, source problem code \texttt{catalogue entry}, collection \emph{Oberwolfach Reports Open Problems}. The supplied record gives proposed year 2019 and names no individual proposer.
\item Supplied source citation: \emph{Modular Forms (2020)}, Oberwolfach Reports, DOI \texttt{10.4171/owr/2019/57}, pp.\ 3529--3606. \url{https://doi.org/10.4171/owr/2019/57}. The supplied background describes Oberwolfach Reports as reports of mathematical workshops and their open-problem discussions; no problem-specific page is supplied.
\item Supplied literature source: Campagna, effective bounds on differences of singular moduli. \url{https://doi.org/10.1017/S0305004122000378}.
\item Supplied literature source: OWR report. \url{https://doi.org/10.4171/owr/2019/57}.
\item Literature assessment checked at \texttt{2026-08-22}: the supplied curation labels the question open and states that fixed-singular-modulus finiteness is known, while the question with both moduli varying remains conjectural. This is supplied dated curation derived from the OWR material, not an independently verified assessment. The supplied record explicitly requests independent verification of time-sensitive status and cited claims.
\end{itemize}
\end{document}
